% style-lint: six-birds-section
\section{The Six Birds framework behind this paper}
\label{app:sixbirds}

This appendix is the only place where we use the vocabulary of the Six Birds
framework, a programme of structural laws for emergence and closure across
layers. The paper's questions and the design of its models come from that
programme, and this appendix says how. None of our results depends on any
statement below. Each result stands on its own proof in Lean.

\paragraph{The question in the framework's terms.}
In the framework, a layer is promoted when a package of conditions closes on a
carrier, and the consciousness project specifies such a package, the
C-package with clauses C1 to C8 \citep{Tsiokos2026AuditedClosureArchitecture,
TsiokosArchitecture2026}. The project behind this paper, there called
substrate transcendence, asks when a promoted package can be carried from one
carrier to another. It separates the three lanes the body calls realization,
lineage and autonomy, and prices each separately. The strict criterion of the
body is the framework's Temporal-FIACH certificate. Continuation by ancestry is
its Temporal-IAPS fallback.

\paragraph{The sources the design drew on.}
The design of the models drew on statement-level sources in the framework's
corpus:
\begin{itemize}
\item the C-package contract and its temporal clause
\citep{Tsiokos2026AuditedClosureArchitecture};
\item predictive equivalence, predictive transport and route residue on fixed
support \citep{Tsiokos2026HolonomyMemory};
\item branchwise versus recombination structure
\citep{Tsiokos2026MarkingErasure};
\item the catalog laws for holonomy-memory repair, relation-valued object
persistence and arrows from record monotonicity \citep{Tsiokos2026FoundationsIV};
\item holonomy without an endogenous arrow \citep{Tsiokos2026ProtocolTrap};
\item flatness, local confluence and unique normal forms
\citep{Tsiokos2026FlattenStone};
\item time as order, measure and arrow, funded clocks, and clock translation
\citep{Tsiokos2026NotchStone};
\item viability and target-owned maintenance \citep{Tsiokos2026Agenthood};
\item an analogy with logical objects in quantum error correction
\citep{Tsiokos2026QuantumErrorCorrection}.
\end{itemize}
Each source carries its own scope warning. Most results are on fixed support
and do not by themselves give a cross-carrier theorem. The catalog laws are
conditional on declared quotients and instruments. The quantum analogy stays
an analogy.

\paragraph{Versions.}
The design read these sources in the versions frozen on 2026-09-28, which we
cite by content hash in the artifact. Several were revised afterwards. Ten of
the nineteen cited statements have since changed, and three of the changes
are material. The time paper now defines the arrow by path reversal and calls
a stationary monotone ledger trivial. The agency paper no longer credits zero
defect to repair. The holonomy paper no longer states that current transport
exists exactly when the interface is flat. These revisions do not affect our
results, which do not depend on the sources. They do explain one choice of
wording in the body: the record arrow is an ordering by record growth, not an
arrow in the revised time paper's sense.

\paragraph{What the framework does not supply.}
The corpus is not evidence that any system is conscious or that consciousness
can move between substrates. Substrate Indifference, A2 in the architecture,
remains a postulate with no transport or invariance theorem. Certificate
soundness is open, and the architecture's integrated instance has not been
exhibited.
